\documentclass[11pt]{article}

\usepackage[T1]{fontenc}
\usepackage{lmodern}
\usepackage[margin=1in]{geometry}
\usepackage{booktabs}
\usepackage{amsmath,amssymb}
\usepackage{hyperref}
\usepackage{xcolor}
\usepackage{listings}

\newcommand{\toolname}{\textsc{arx}}
\newcommand{\code}[1]{{\small\texttt{#1}}}

\title{\toolname{}: A Statistical Gate for Autonomous Research Loops\\
       \large (What happened when a local LLM wrote 19 experiments)}
\author{}
\date{}

\begin{document}
\maketitle

\begin{abstract}
Autonomous research loops that combine an LLM with tool access are usually
evaluated on how much work they complete. We ask a different question: of the
claims such a loop produces, how many survive a check on whether the
measurement could have resolved them at all?

We built \toolname{}, a loop whose reports are gated by classical measurement
theory rather than by the loop's own judgement. The gate refuses a claim when
the gold producers are also the contestants, when the floor and the board are
computed on different item sets, when fewer than three gold producers exist,
when the reported margin is below the minimum detectable effect, or when the
measured floor is not detectably above the chance rate.

We then ran the complete loop five times on five topics, with a local
7B-parameter model writing both the hypotheses and the experiment code, over
107 papers retrieved from arXiv. The loop produced 19 gate decisions. Two were
statistical verdicts: one floor was not estimable under the candidate
processor's design, and one ``leader'' scored below the measured noise floor.
The other seventeen were failures of code generation, not of statistics---the
gate relayed 13 crashes and 4 outputs that never emitted the contract. It
wrote zero findings.

The headline number, 0/19, is not the contribution. The contribution is the
separation: a loop that reports ``19 hypotheses tested, all inconclusive'' is
describing 17 broken experiments and 2 genuine statistical refusals as if
they were the same thing, and a reader cannot tell which from that sentence.
We describe the gate, the four refusals, the ledger that makes the run
auditable, and what this run does \emph{not} establish.
\end{abstract}

%==========================================================================
\section{Introduction}
%==========================================================================

The past year produced many demonstrations of an LLM driving a research loop:
propose a hypothesis, write code, run it, read the result, repeat. The
evaluation question has almost always been throughput---how many experiments
ran, how many papers were read, whether a report appeared at the end. Almost
never: how many of the resulting claims were \emph{measurements} rather than
plausible-looking text.

That question matters more than it first appears, because the failure is not
random. An LLM asked to produce an experiment produces text that looks like an
experiment: random score tables labeled ``system-0'' through ``system-2'',
majority votes defined with a two-element judge set, agreements computed
against \code{scores[0]} when \code{scores} is a dictionary. It is not lying.
It has no model of what the numbers are for. A loop that pipes that output
into a report converts fluent code into a scientific claim, and the report
gives no signal which one happened.

\toolname{} is one response: put a gate between the run and the report, where
the gate is classical measurement theory rather than a second LLM asked to be
critical. The gate cannot fix a broken experiment. It can refuse to describe
one as a result.

\paragraph{What this paper is not.} We are not claiming a new statistical
method. Leave-one-out floors, the four preconditions, and the paired
minimum-detectable-effect calculation are all textbook. We are not claiming
the loop is useful for research: five topics and 19 experiments is a
mechanism of evaluation, not a research result. And we are not claiming the
gate works well against a smarter model---it has been run against one 7B
model and nothing else.

%==========================================================================
\section{The gate}
%==========================================================================

A leaderboard is \emph{readable} when its ordering separates the systems more
than the measurement apparatus separates itself from nothing. Four
preconditions, each a reason to refuse:

\begin{enumerate}
\item \textbf{Actor-disjointness.} The systems that produced the reference
labels are not the systems being ranked. If the gold producers are also the
contestants, the floor partly measures agreement with itself, and a contestant
that helped define the gold has an advantage no amount of averaging removes.
\item \textbf{Common denominator.} The floor and the board are computed on the
same item set. The difference between a floor on 1,520 items and a board on
1,084 of them is not an effect size; it is an artefact of subsetting, and
subtracting the two produces a number that means nothing.
\item \textbf{Multi-generator sufficiency.} At least three reference-label
producers. With two, removing one leaves one, so ``agreement with the
others'' is agreement with that one. The quantity is not weak; it does not
exist.
\item \textbf{Power-bounded resolution.} The margin between leader and floor
exceeds the minimum detectable effect at the declared power. A margin of
$+0.03$ with an MDE of $0.05$ is not a small effect, it is a sample too small
to have an opinion.
\end{enumerate}

A fifth check runs first and quietly decides most cases. The floor is only
interpretable if it is detectably above the chance rate. With $K$ labels the
chance level is $1/K$; with five generators each flipping a coin, pairwise
agreement sits near $1/K$ and every score above it looks like a separation. So
the floor is tested against chance with the standard error that comes from the
item count,
\[
z \;=\; \frac{\text{floor} - 1/K}{\sqrt{(1/K)(1-1/K)/n}}\,,
\]
and a $z$ below $z_{1-\alpha/2}$ means the reference labels have not
demonstrably reproduced anything, and there is no floor to compare a score
against.

\paragraph{Why the ceiling is checked before the wall.} A board whose floor is
at chance with a margin of $+0.47$ is the loudest possible false positive, and
it survives every margin-only check: the margin is large and positive. Ordering
the chance test first is the single highest-value decision in the gate.

%==========================================================================
\section{What the gate cannot do}
%==========================================================================

\begin{itemize}
\item It cannot repair an experiment. A crashed run produces no verdict at
all, only a recorded non-zero exit code.
\item It cannot tell a real effect from an artefact. \code{READABLE} says the
apparatus can separate this claim from its own noise; it says nothing about
whether the apparatus measured the right thing.
\item It cannot stop a loop from producing more hypotheses after a refusal.
What it does is refuse to write them up, which is the only part a downstream
reader sees.
\end{itemize}

%==========================================================================
\section{Method: what actually ran}
%==========================================================================

Five runs on five topics, each starting from a single natural-language
question, with a local \code{llama3.2:3b} (Ollama, OpenAI-compatible endpoint)
writing both the hypotheses and the experiment scripts. Literature came from
the arXiv API; failures to reach a source were recorded, not hidden.

The loop, at each step, chooses among five actions by a hand-written
information-gain heuristic divided by a cost, where the cost of re-running the
literature search while candidates are untested is set high on purpose. It
terminates on a stall counter or a wall-clock ceiling. Every decision, the
budget snapshot, and the gate's reasons are appended to a hash-chained ledger.

The experiment contract is a single JSON object on stdout:

\lstset{basicstyle=\footnotesize\ttfamily,frame=single}
\begin{lstlisting}
{"votes": [[0,1,-1], ...],        // m generators x n items, -1 abstain
 "scores": {"system-0": 0.91},   // contestant accuracies
 "judge_ids": ["g1","g2","g3"],
 "contestant_ids": ["system-0", ...],
 "item_ids": [0,1,...],           // denominator, for the mismatch check
 "se": 0.02}
\end{lstlisting}

\paragraph{Pre-registration as ordering, not policy.} Before the subprocess
starts, the loop decides the number of planned looks, hashes the script it is
about to execute, and writes both to the ledger. The gate then reads
\code{k\_max} from the ledger, never from the experiment's own report, and
compares the hash measured off the file that ran against the one pinned. A
script that reports a different look count or a different hash is refused
before its numbers are examined. This closes the ``run a fixed-sample test,
miss, then claim you only intended one look'' hole structurally: the claim
becomes impossible rather than forbidden.

%==========================================================================
\section{Results}
%==========================================================================

\begin{table}[t]
\centering\small
\begin{tabular}{lrrrrl}
\toprule
topic & papers & exps & crashes & \multicolumn{2}{c}{gate verdicts}\\
\midrule
reward hacking                     & 20 & 3 & 2 & \multicolumn{2}{l}{$3\times$ UNDECIDABLE}\\
noise floors in LLM-as-judge       & 27 & 5 & 3 & \multicolumn{2}{l}{$5\times$ UNDECIDABLE}\\
calibration of confidence          & 20 & 4 & 1 & \multicolumn{2}{l}{$1\times$ UNREADABLE, $3\times$ UNDECIDABLE}\\
spurious correlations              & 20 & 3 & 3 & \multicolumn{2}{l}{$3\times$ UNDECIDABLE}\\
reproducibility of margins         & 20 & 4 & 4 & \multicolumn{2}{l}{$4\times$ UNDECIDABLE}\\
\midrule
\textbf{total} & \textbf{107} & \textbf{19} & \textbf{13} &
\end{tabular}
\caption{Five complete loop runs. ``Crashes'' counts experiments that exited
non-zero. All five ledgers verify as valid hash chains.}
\label{tab:runs}
\end{table}

\subsection{The split that matters}

The 19 gate decisions are not 19 of the same thing:

\begin{table}[t]
\centering\small
\begin{tabular}{lr}
\toprule
reason for the gate's decision & count\\
\midrule
execution failure the gate merely relayed & \textbf{17}\\
\quad experiment exited non-zero & 13\\
\quad no JSON object found on stdout & 4\\
\addlinespace
a verdict the gate reached statistically & \textbf{2}\\
\quad floor not estimable (\code{floor\_degenerate}) & 1\\
\quad leader below the measured floor (\code{UNREADABLE}) & 1\\
\midrule
findings written & \textbf{0}\\
\bottomrule
\end{tabular}
\caption{The same 19 decisions, separated by whether the gate did any
statistical work. Reporting only ``0 findings'' or only ``19 decisions''
loses this distinction.}
\label{tab:split}
\end{table}

\subsection{The two statistical refusals}

\paragraph{The floor did not exist.} One run produced five generators over ten
items and the gate returned \code{floor\_degenerate}. This is the case where
the candidate processor's own design was insufficient, and the gate's verdict
was the honest one. No amount of analysis of that output produces a floor, and
the alternative---computing something and labelling it a floor---is the failure
this gate exists to prevent.

\paragraph{The ``leader'' was below the noise.} One run scored a contestant at
0.80 against a floor of 1.00 (ten items, three generators). The margin was
$-0.20$. This is the case where a loop without a gate would report ``the
leader scores 0.80'' as a finding, and the number 0.80 is arithmetically true
and scientifically empty.

\subsection{What the 13 crashes looked like}

Not adversarial inputs. Representative failures from the model's own
experiment code:

\begin{itemize}
\item \code{NameError: name 'i' is not defined}---a loop variable used
outside the loop that introduced it.
\item \code{KeyError: 0} on \code{scores[vote[0]]}, where \code{scores} is a
dictionary keyed by strings and \code{vote[0]} is an integer index. The code
had no model of which it was.
\item \code{Object of type float32 is not JSON serializable}---output produced
by \code{json.dumps} on numpy scalars.
\end{itemize}

These are not prompt failures. They are the ordinary consequence of asking a
language model for ``an experiment'' without giving it a model of what the
numbers are for, and they are the reason a gate has to sit in the path rather
than an instruction that says ``do not fabricate results.''

%==========================================================================
\section{The honest reading of 0 findings}
%==========================================================================

The number 0/19 answers one question: \emph{did any claim produced by this
loop clear a statistical check on its measurement?} It does not answer whether
a stronger model would produce experiments that run, whether the gate would
then let bad claims through, or whether the topics were ones where a small
experiment could learn anything.

It also does not show the gate's selectivity, because selectivity requires
candidates that can fail on statistical grounds rather than execution grounds.
With 2 statistical decisions, the gate has been exercised, not validated. That
is a fact about this run and we report it rather than rounding it up.

What the run does establish is narrower and more useful: when the loop is
allowed to write its own report, all 19 decisions land in the report with no
way for a reader to tell 17 broken experiments from 2 refusals. With the gate
and the ledger, the split in Table~\ref{tab:split} is read back mechanically,
and the ledger hash chain means the split can be checked by someone who was
not there.

%==========================================================================
\section{Related work}
%==========================================================================

LLM-as-a-judge evaluation and its noise floors are covered in the concurrent
literature on judge reliability; our floor estimate is the leave-one-out
agreement used in that line of work. Agent frameworks for science typically
end at producing a report; \toolname{} differs in that the report is not
writable without a passing audit.

%==========================================================================
\section{Limitations}
%==========================================================================

\begin{itemize}
\item One model, one size, five topics. Nothing here is a benchmark and we do
not claim a number that would generalize.
\item 68\% of experiments crashed, so the gate was exercised on 2 statistical
vent points. Selectivity is untested.
\item The information-gain heuristic that orders the loop's actions is
hand-written, not fitted. It is coarse and it is not the subject of this
paper; a fitted version would immediately become a target of the same
question this paper asks, which is why it is deliberately not learned here.
\end{itemize}

%==========================================================================
\section{Conclusion}
%==========================================================================

The interesting number in this run is not 0/19 findings. It is 17/19 execution
failures and 2/19 statistical refusals, and the fact that without a gate those
19 arrive in a report looking identical.

\end{document}
